\documentclass[10pt]{article}
\usepackage[margin=1in]{geometry}
\usepackage[T1]{fontenc}
\usepackage{lmodern,microtype,amsmath,amssymb,amsthm}
\theoremstyle{definition}
\newtheorem{assumption}{Assumption}
\begin{document}
\section*{Momentum SGD for Convex Minimization}
We consider minimizing a convex function $f:E\to\mathbb{R}$ on a separable real Hilbert space $E$. A given point $x^0\in E$ lies within distance $R\geq0$ of a global minimizer $x^\star$. At $x$, the oracle returns $g(\xi,x)\in E$ using a fresh independent sample $\xi$ from a probability law $\nu$ on a measurable space $\Xi$. The map $g$ is jointly measurable, its mean is a subgradient of $f$, and $\mathbb{E}_{\xi\sim\nu}\|g(\xi,x)\|^2\leq G^2$, where $G>0$. For $\varepsilon>0$, constant-momentum SGD must return a last iterate $\bar{x}$ with integrable objective gap satisfying $\mathbb{E}[f(\bar{x})-f(x^\star)]\leq\varepsilon$ under the sampling law.

The optimization problem is
\[
\min_{x\in E} f(x).
\]
The norm is induced by the inner product on $E$. The algorithm receives $E$, the measurable sample space $\Xi$, $G$, $R$, a momentum parameter $\beta\in[0,1)$, $\varepsilon$, and $x^0$; the objective, oracle map, and sampling law are accessed through oracle replies.

Our assumptions:
\begin{assumption}[Convexity]
For every $x,y\in E$ and $t\in[0,1]$,
\[
f((1-t)x+ty)\leq(1-t)f(x)+tf(y).
\]
\end{assumption}

\begin{assumption}[Stochastic subgradient oracle]
The map $g:\Xi\times E\to E$ is jointly measurable. For every $x\in E$, the random variable $\|g(\xi,x)\|^2$ is integrable under $\nu$, and
\[
\mathbb{E}_{\xi\sim\nu}\bigl[\|g(\xi,x)\|^2\bigr]\leq G^2.
\]
Writing $h(x)=\mathbb{E}_{\xi\sim\nu}[g(\xi,x)]$ for the Bochner mean, we require
\[
f(x)+\langle h(x),y-x\rangle\leq f(y),
\qquad x,y\in E.
\]
Successive calls use independent samples $\xi_0,\xi_1,\ldots$ with common law $\nu$, and queries may be made anywhere in $E$.
\end{assumption}

\begin{assumption}[Initial distance]
There exists $x^\star\in E$ such that
\[
\begin{aligned}
f(x^\star)&\leq f(x) &&\text{for every }x\in E,\\
\|x^0-x^\star\|&\leq R.
\end{aligned}
\]
\end{assumption}

Our result is algorithm-specific: we analyze constant-momentum SGD with a linearly decreasing step-size schedule. The algorithm chooses $\gamma\in\mathbb{R}$ and $N\in\{0,1,2,\ldots\}$ from its inputs and initializes $m^0=0$. At iteration $k\geq0$, it queries the oracle at $x^k$ and updates
\[
\begin{aligned}
g_k&=g(\xi_k,x^k),\\
m^{k+1}&=\beta m^k+(1-\beta)g_k,\\
\eta_k&=\gamma\left(N+1-k+\frac{\beta}{1-\beta}\right),\\
x^{k+1}&=x^k-\eta_k m^{k+1}.
\end{aligned}
\]
After $T\in\{0,1,2,\ldots\}$ oracle calls, the output is the last iterate $\bar{x}=x^T$. The schedule parameter $N$ and the call budget $T$ are separate quantities.

Let $P=\nu^{\otimes\mathbb{N}}$ be the joint law of the samples. The required objective gap is integrable under $P$ and satisfies
\[
\mathbb{E}_{P}\bigl[f(\bar{x})-f(x^\star)\bigr]\leq\varepsilon.
\]
\end{document}
